\section{Localization and joint spatial estimates}\label{sec:spatial}

We now show that any prescribed family of cells is unlikely to suffer
simultaneous particle shortages or simultaneous failures of separation.
These are estimates under the actual equilibrium position law $\mu$.
They do not condition on the positions outside the cells. The method,
adapted from \citet[Section~4]{B}, compares a lower free-energy penalty
from spatial localization with an upper bound obtained by tilting the
diagonal of the Gibbs state.

For a bounded measurable one-particle function $f$, put
$n_f(X)=\sum_i f(x_i)$, and write $\E_\sigma$ for expectation under
the position law of a state $\sigma$. In this section set
\[
 \Gamma=Z_N^{-1}e^{-\beta H_N},\qquad
 \mu(\dd X)=\Gamma(X,X)\dd X.
\]

\subsection{Localization and diagonal tilting}

\begin{lemma}[Two-color localization]\label{lem:localization}
Let $p,q$ be real Lipschitz functions with $p^2+q^2=1$, and suppose
$p$ is compactly supported in an open set $O$, or take $O=\Lambda_K$.
For a finite-energy symmetric $J$-particle state $\sigma$, map each
coordinate by $f\mapsto(pf,qf)$ and dephase by the number $j$ of
particles in the first color. Denote its distribution by $w_j$ and
the conditional first- and second-color marginal states by
$\sigma_{i,j}$ and $\sigma_{o,j}$. For every bounded field $W\ge0$,
\begin{equation}\label{eq:two-color}
 \begin{split}
 \mathcal F_J^{(-W)}(\sigma)
  +\E_\sigma n_{|\nabla p|^2+|\nabla q|^2}
 \ge\sum_{j=0}^Jw_j\bigl[
  \mathcal F_j^{(-W)}(\sigma_{i,j})+
  \mathcal F_{J-j}^{(-W)}(\sigma_{o,j})\bigr]-TH(w).
 \end{split}
\end{equation}
The inside states also satisfy the Dirichlet condition in $O$. Terms
with $w_j=0$ are omitted, and vacuum states are allowed. Moreover
\begin{equation}\label{eq:color-entropy}
 \sum_jjw_j=\E_\sigma n_{p^2},\qquad
 H(w)\le-\log(1-e^{-1})+\sum_jjw_j.
\end{equation}
\end{lemma}

\begin{proof}
The geometric localization and entropy decomposition are the many-body
IMS construction; see \citet[Section~3]{Lewin2011} and
\citet[Section~3, Lemma~3.1]{DeuchertSeiringerYngvason2019}.
Here is the hard-core version. The one-particle map is isometric, so
its symmetric tensor power preserves the spectrum and entropy of
$\sigma$. In each sector of fixed color number, normalized
symmetrization over the mutually orthogonal color assignments identifies
the space with the tensor product of the two symmetric species spaces.
There is one such tensor product, with no binomial multiplicity.

The product rule and $p\nabla p+q\nabla q=0$ give the exact
kinetic increase $\E_\sigma n_{|\nabla p|^2+|\nabla q|^2}$.
The field distributes between the two colors without error. The masks
preserve the hard-core form domain. After cross-color exclusions are
dropped, slicing gives admissible same-color marginal states, by the
domain argument in Section~\ref{sec:energy}; the first color is also
Dirichlet in $O$. Thus the sum of their energy expectations equals
the localized energy. Entropy cannot decrease under dephasing and is
subadditive in each color sector, whence
\[
 S(\sigma)\le H(w)+\sum_jw_j
       [S(\sigma_{i,j})+S(\sigma_{o,j})].
\]
All entropies are finite by Lemma~\ref{lem:entropy}, applied also to
the finitely many color spaces. This proves \eqref{eq:two-color}.
Conditional on positions, colors are independent choices with
probabilities $p(x_i)^2,q(x_i)^2$, proving the first identity in
\eqref{eq:color-entropy}. Nonnegativity of relative entropy to the
geometric law $(1-e^{-1})e^{-j}$ proves the second.
\end{proof}

\begin{lemma}[Diagonal tilt]\label{lem:tilt}
For a bounded real symmetric Lipschitz function $G$ of $N$ positions,
set
\[
 g=e^G,\qquad M_G=\E_\mu e^{2G},\qquad
 \sigma_g=M_G^{-1}g\Gamma g.
\]
Then
\begin{equation}\label{eq:diagonal-tilt}
 \mathcal F_N(\sigma_g)-F_N
 \le\E_{\sigma_g}|\nabla G|^2
       +T\bigl(2\E_{\sigma_g}G-\log M_G\bigr).
\end{equation}
Here $|\nabla G|^2=\sum_i|\nabla_iG|^2$.
\end{lemma}

\begin{proof}
This is the diagonal-tilting argument of
\citet[Section~4, ``Two-color localization and diagonal tilting'']{B}.
Put $\tau=M_G^{-1}\Gamma^{1/2}g^2\Gamma^{1/2}$.
The nonzero spectra of $\tau$ and $\sigma_g$ agree, as the
spectra of $B^*B$ and $BB^*$ for
$B=M_G^{-1/2}g\Gamma^{1/2}$. If $H_N\psi=E\psi$,
testing the weak equation with $g^2\psi$ gives
\[
 q_N[g\psi]=E\norm{g\psi}_2^2+
                      \int|\nabla g|^2|\psi|^2.
\]
Lipschitz multiplication preserves the hard-core boundary condition.
Summing in a Gibbs eigenbasis yields
\begin{equation}\label{eq:tilt-energy-identity}
 \Tr H_N\sigma_g=\Tr H_N\tau+
                         \E_{\sigma_g}|\nabla G|^2.
\end{equation}
Both energies are finite: $g$ and its gradients are bounded at fixed
$N,K$, and positive-time heat traces control the eigenvalue-weighted
sum. Lemma~\ref{lem:entropy} then gives finite entropies.

To bound the relative entropy $D(\tau\Vert\Gamma)$, first let
$g^2$ be constant on a finite partition into symmetric measurable
sets $E$. Put
$p_E=\Tr\Gamma\ind_E$ and
$\tau_E=p_E^{-1}\Gamma^{1/2}\ind_E\Gamma^{1/2}$
when $p_E>0$. The two states $\Gamma$ and $\tau$ are mixtures
of the same states $\tau_E$, with weights $p_E$ and
$p_Eg_E^2/M_G$, respectively. Monotonicity of relative entropy under
discarding the classical label \citep{Lindblad1975} gives
\[
 D(\tau\Vert\Gamma)\le
 \sum_E\frac{p_Eg_E^2}{M_G}\log\frac{g_E^2}{M_G}.
\]
Uniform positive simple approximation to $g^2$ gives trace-norm
convergence of these states and convergence of the classical expression.
Lower semicontinuity therefore implies
$D(\tau\Vert\Gamma)\le2\E_{\sigma_g}G-\log M_G$.
Finally
$\mathcal F_N(\tau)-F_N=TD(\tau\Vert\Gamma)$,
the equality of entropies, and \eqref{eq:tilt-energy-identity} prove
\eqref{eq:diagonal-tilt}.
\end{proof}

\subsection{A local penalty and density control in a field}

For a nonempty set of sites $A\subset\mathcal L_K$, define, for a
configuration of any particle number $J\le N$,
\begin{align}
 s_0&=bD^{-1/3},\qquad 0<b<1,\label{eq:s0}\\
 f_A(X)&=\sum_{i=1}^J e^{-d(x_i,A)},\label{eq:weighted-count}\\
 k_A(X)&=\#\{i:d(x_i,A)\le6,\ d_i\le2s_0\}.
                                      \label{eq:localized-close-count}
\end{align}
The parameter $b$ will be chosen small after the low-count threshold
has been fixed. We always increase $D$ so that $r<s_0<1/2$.

\begin{proposition}[Weighted local free-energy penalty]\label{prop:local-penalty}
There are $c,C>0$, independent of $b$, such that for every
$0\le J\le N$, every bounded measurable $0\le W\le1$, and every
finite-energy symmetric trace-one state $\sigma$,
\begin{equation}\label{eq:local-penalty}
 \mathcal F_J^{(-W)}(\sigma)-F_J^{(-W)}
 \ge\alpha\left[\E_\sigma(f_A+cb^{-3}k_A)-CD|A|\right].
\end{equation}
The threshold on $D$ may depend on $b$, but the displayed constants
do not.
\end{proposition}

\begin{proof}
We combine the exponential localization of
\citet[Section~3, ``Exponential spatial localization'']{HS} with
the thermal entropy decomposition of Lemma~\ref{lem:localization}.
For $h\ge0$, choose sine--cosine masks $p_h,q_h$ such that
$p_h=1$ at distance at most $h$ from $A$, $p_h=0$ from distance
$h+1$ onward, and
\[
 |\nabla p_h|^2+|\nabla q_h|^2
 \le C\ind_{\{h<d(x,A)<h+1\}}.
\]
Take the inside region $O_h$ to be the open $(h+2)$-neighborhood,
or the torus when appropriate. Its volume is at most
$C(h+4)^3|A|$, and $p_h$ is compactly supported there.

Apply Lemma~\ref{lem:localization}. Outside, the upper increments of
Proposition~\ref{prop:increments} give
$F_{J-j}^{(-W)}\ge F_J^{(-W)}-C\alpha j$.
Inside, divide the kinetic energy into three fixed positive portions.
One portion pays for the entire von Neumann entropy by
Lemma~\ref{lem:entropy}. Apply Lemmas~\ref{lem:cube} and
\ref{lem:close} to the other two, with $s=s_0$. The field costs
at most $j$, while the color entropy costs at most $T(C+j)$ by
\eqref{eq:color-entropy}. Thus all losses linear in the inside
number are bounded by $C\alpha j$, since $T\le1$ and $\alpha\ge1$.
The remaining entropy error is at most
\[
 CT\bigl[1+|O_h|(V^{-1}+T^{3/2})\bigr]
 \le C(h+4)^3|A|.
\]

Average with the probability density $e^{-h}\,\dd h$.
For every position $x$,
\[
 \int_0^\infty p_h(x)^2e^{-h}\,\dd h\le e\,e^{-d(x,A)},
 \quad
 \int_0^\infty\ind_{\{h<d(x,A)<h+1\}}e^{-h}\,\dd h
 \le e\,e^{-d(x,A)}.
\]
Consequently the averaged removal and localization costs are at most
$C\alpha\E_\sigma f_A$, and the residual entropy cost is at most
$C\alpha|A|$.

To retain the positive terms, set $b_v=e^{-d(v,A)}$. If
$h\ge d(v,A)+1$, every coordinate in $B_v$ belongs to the first
color with probability one. The averaged cube contribution is therefore
at least
$cr\E_\sigma\sum_v b_vn_v(n_v-1)$.
A particle counted by $k_A$ has a partner within $2s_0<1$;
both lie at distance less than $7$ from $A$. For $h\ge8$ both
labels belong to the first color, so its close-neighbor inequality
counts this particle. This yields at least
$crs_0^{-3}\E_\sigma k_A=c\alpha b^{-3}\E_\sigma k_A$.
The color interpretation used here follows from the conditional
independent choices in Lemma~\ref{lem:localization}.

We have proved
\[
 \mathcal F_J^{(-W)}(\sigma)-F_J^{(-W)}
 \ge\E_\sigma\left[cr\sum_vb_vn_v(n_v-1)
       -C\alpha f_A+c\alpha b^{-3}k_A\right]-C\alpha|A|.
\]
Within a unit cell, distances to $A$ differ from the center distance
by less than one. Hence
\[
 f_A\le e\sum_v b_vn_v,\qquad
 \sum_v b_v\le\sum_{w\in A}\sum_v e^{-d(v,w)}\le C|A|.
\]
The final lattice sum is uniformly bounded on periodic unit grids.
For $n\ge0$ and $D\ge1$, completing a scalar square gives
\[
 \frac cD n(n-1)-(C+1)e n\ge-C'D.
\]
Multiply by $\alpha b_v$ and sum. Since $r=\alpha/D$, this
absorbs all linear losses while leaving $\alpha f_A$ and proves
\eqref{eq:local-penalty}. The vacuum case is immediate.
\end{proof}

\begin{corollary}[Density bound in a small field]\label{cor:field-density}
In the $J$-particle equilibrium state with $0\le J\le N$ and
$0\le W\le1$, every set $A$ of sites satisfies
\begin{equation}\label{eq:field-density}
 \E n_{\ind_{\bigcup_{v\in A}B_v}}\le CD|A|.
\end{equation}
\end{corollary}

\begin{proof}
The left side of \eqref{eq:local-penalty} is zero in equilibrium.
Drop the nonnegative close-count term and use
$n_{\ind_{\cup_{v\in A}B_v}}\le e f_A$. Empty $A$ and the
vacuum cause no exception.
\end{proof}

The mean estimate will control the effect of small fields on partition
functions. To produce many movable particles, we need the stronger joint
statements proved next.

\subsection{Simultaneous shortages are exponentially unlikely}

\begin{proposition}[Joint low-count bound]\label{prop:low-count}
One can choose $\alpha_0$ sufficiently large and then $\ell>0$
sufficiently small so that, for every set of sites $A$ and all
sufficiently large $D$,
\begin{equation}\label{eq:low-count}
 \mu\{n_v<\ell D\text{ for every }v\in A\}
 \le e^{-cD|A|}.
\end{equation}
The constants and thresholds are uniform in $K\ge100$ and
$0<T\le D^{-1}$.
\end{proposition}

\begin{proof}
Assume $A\ne\varnothing$. In each of its cells take $p=1$ on
the central cube of half-side $1/4$, with $p=0$ outside the central
cube of half-side $1/3$. Sine--cosine interpolation gives $q$ with
$p^2+q^2=1$ and uniformly bounded squared gradient cost. Let $O$
be the union of these unit-cell interiors. Choose a smooth squared
cutoff $0\le\phi\le1$, supported in $O$, equal to one on the
central cubes of half-side $1/3$, and satisfying
$|\nabla\phi|^2\le C\phi$. In particular $p^2$ and the
localization gradient cost are bounded by $C\phi$.

We compare each outside state from Lemma~\ref{lem:localization}
with a state obtained by filling the empty central cubes. Set
$k=\lfloor\eta D|A|\rfloor$, where $0<\eta<1$ will be fixed.
Take a normalized smooth bump in each central cube of half-side
$1/8$, and sum with equal amplitude to obtain an orbital $w$ with
\[
 \norm w_\infty^2\le C/|A|,\qquad
 \norm{\nabla w}_2^2\le C.
\]
Starting at the vacuum, apply Lemma~\ref{lem:insertion}
successively $k$ times with this orbital. Its smallness condition is
satisfied because
$Ckr^3/|A|\le CD r^3\to0$. The resulting normalized
labelled vector has energy at most
\begin{equation}\label{eq:fill-cost}
 Ck(1+rk/|A|).
\end{equation}
It is supported in the union of the bump supports in every coordinate.
Replace it by the square root of the average of its squared moduli
over all permutations. The norm is unchanged, the hard-core zeros
are preserved, and the inequality for the gradient of a Euclidean
norm shows that energy cannot increase. This constructs a fixed
symmetric normalized added vector $\psi_k$ with the same bound.
Dirichlet admissibility follows by contraction and approximation of
the original form vector.

The conditional outside state is supported in
$(\supp q)^{N-j}$. Its coordinates are separated from every bump
support by distance at least $1/8$, larger than $r$ for large
$D$. The map on its form vectors
\[
 U_jf=\binom{N-j+k}{k}^{1/2}
       P_{\rm sym}(f\otimes\psi_k)
\]
is an isometry: the distinct choices of the $k$ added labels have
disjoint position supports. It preserves entropy and has additive
kinetic energy. The support separation also guarantees every
old--new hard-core condition. Consequently
\[
 \mathcal F_{N-j}(\sigma_{o,j})
 \ge F_{N-j+k}-Ck(1+rk/|A|).
\]
Here $k\le\eta N<N$, so Proposition~\ref{prop:increments}
applies up to $N+k\le2N$. Its lower increment followed by $j$
upper increments gives
\begin{equation}\label{eq:outside-fill-comparison}
 \mathcal F_{N-j}(\sigma_{o,j})
 \ge F_N+c\alpha k-C\alpha j-Ck(1+rk/|A|).
\end{equation}
The bump profiles and their constants were fixed independently of
$\alpha_0$. Choose $\alpha_0$ large, and then $\eta$ small,
so that throughout the $\alpha$ window
\[
 C(1+\alpha\eta)\le c\alpha/2.
\]
Since $rk/|A|\le\alpha\eta$, the final term in
\eqref{eq:outside-fill-comparison} is absorbed by half the positive
insertion gain. For large $D$, $k\ge\eta D|A|/2$.

Inside, Lemma~\ref{lem:entropy} with the full kinetic energy bounds
the free energy below by $-Tj-CT|A|(V^{-1}+T^{3/2})$.
Apply Lemma~\ref{lem:localization}, use
$\sum_jjw_j\le\E_\sigma n_\phi$, and absorb the remaining
$O(|A|)$ errors into the filling gain. For every finite-energy
$N$-particle state this gives
\begin{equation}\label{eq:shortage-penalty}
 \mathcal F_N(\sigma)-F_N
 \ge c\alpha D|A|-C\alpha\E_\sigma n_\phi.
\end{equation}

Apply Lemma~\ref{lem:tilt} with $G=-\lambda n_\phi$,
$0\le\lambda\le1$, and write
$L(\lambda)=-\log\E_\mu e^{-2\lambda n_\phi}$.
Since $0\le n_\phi\le N$, this moment is differentiable at fixed
$N,K$ by dominated convergence.
The squared gradient is at most $C\lambda^2n_\phi$, while
$2\E_{\sigma_g}G\le0$. Comparing with
\eqref{eq:shortage-penalty}, and increasing $C$, gives
\[
 C\alpha\E_{\sigma_g}n_\phi+TL(\lambda)
       \ge c\alpha D|A|,\qquad
 L'(\lambda)=2\E_{\sigma_g}n_\phi.
\]
Choose a fixed $a'>0$ sufficiently small. Whenever
$L(\lambda)\le a'D|A|$, the term $TL(\lambda)$ is at most
half the right side, and the resulting lower bound for $L'$ is at
least $2a'D|A|$. Since $L(0)=0$ and $L$ is nondecreasing, this
forces $L(1)\ge a'D|A|$. On the low-count event,
$n_\phi\le\sum_{v\in A}n_v<\ell D|A|$. Exponential Markov
therefore gives probability at most
$\exp[-(a'-2\ell)D|A|]$. Choose $2\ell<a'$ to prove
\eqref{eq:low-count}. For empty $A$ both sides equal one.
\end{proof}

\subsection{Simultaneous failures of separation}

Fix from now on $\theta>0$ so that $3\theta<\ell/2$. For a site
$v$, define
\begin{equation}\label{eq:close-count-definition}
 C_v(X)=\#\{i:d(x_i,v)\le5,\ d_i\le s_0\}.
\end{equation}
The enlarged radius allows this count to control endpoint defects of
particles starting in neighboring cells.

\begin{proposition}[Joint close-count bound]\label{prop:close-count}
For every fixed $\theta>0$, one can choose $b>0$ sufficiently small
and then $D$ sufficiently large so that, for every set of sites $A$,
\begin{equation}\label{eq:close-count-tail}
 \mu\{C_v\ge\theta D\text{ for every }v\in A\}
 \le p_D^{|A|},\qquad
 p_D=\exp(-c\theta bD^{2/3})\longrightarrow0.
\end{equation}
Here $c>0$ can be numerical, and the bounds remain uniform in the
temperature and volume ranges above.
\end{proposition}

\begin{proof}
For nonempty $A$, choose squared Lipschitz cutoffs $u,\xi$ in
$[0,1]$, with $u=1$ on $[0,5]$, $u=0$ on $[6,\infty)$,
$\xi=1$ on $[0,1]$, and $\xi=0$ on $[2,\infty)$.
Their derivatives satisfy $|u'|^2\le Cu$ and
$|\xi'|^2\le C\xi$. Define the Lipschitz enlargement
\[
 m_A(X)=\sum_i u(d(x_i,A))\xi(d_i/s_0).
\]
For configurations with fewer than two particles take the summand
with $d_i=+\infty$ to be zero. Then
\begin{equation}\label{eq:close-enlargement}
 m_A\le k_A,\qquad
 \sum_{v\in A}C_v\le C_0m_A,\qquad
 |\nabla m_A|^2\le Cs_0^{-2}m_A.
\end{equation}
The middle inequality follows because each point lies within distance
$5$ of at most a fixed number $C_0$ of unit lattice sites, and each
counted point contributes one to $m_A$.

For completeness, the final inequality uses the geometry of the
nearest-neighbor graph, as in
\citet[Section~3, ``Joint count and close-neighbor tails'']{HS}.
Away from a null set the nearest neighbor of each particle is unique.
If two incoming edges at $x_k$ come from $x_i,x_j$, with
$|x_i-x_k|\ge|x_j-x_k|$, nearest-neighbor choice implies
$|x_i-x_j|\ge|x_i-x_k|$. The cosine law then bounds the angle
between the incoming directions below by $\pi/3$. Edges relevant
to the derivatives have length at most $2s_0<1$, so consistent
Euclidean lifts are available on our torus. Sphere packing therefore
gives a numerical bound on their incoming degree. Each coordinate
derivative meets only a bounded number of summands of $m_A$.
The squared-cutoff inequalities bound the squared full gradient of
each summand by
$Cs_0^{-2}u(d(x_i,A))\xi(d_i/s_0)$.
Bounded overlap proves \eqref{eq:close-enlargement}.

Use Lemma~\ref{lem:tilt} with $G=\lambda m_A$ for
$0\le\lambda\le s_0$. Since $M_G\ge1$,
\eqref{eq:close-enlargement} bounds its right side by
\[
 (C\lambda^2s_0^{-2}+2T\lambda)\E_{\sigma_g}m_A
 \le C\E_{\sigma_g}m_A.
\]
Proposition~\ref{prop:local-penalty}, with $W=0$ and $J=N$,
bounds the same free-energy excess below by
$c\alpha b^{-3}\E_{\sigma_g}m_A-C\alpha D|A|$.
Choose $b$ small enough to absorb the upper bound. It follows that
\[
 \E_{\sigma_g}m_A\le Cb^3D|A|
       \qquad(0\le\lambda\le s_0).
\]
Since $0\le m_A\le N$, differentiating under the integral and then
integrating the logarithmic moment derivative gives
\[
 \log\E_\mu e^{2s_0m_A}\le Cs_0b^3D|A|.
\]
On the event in \eqref{eq:close-count-tail},
$m_A\ge\theta D|A|/C_0$. Reduce $b$ further so that
$Cb^3\le\theta/C_0$. Exponential Markov bounds its probability
by $\exp[-(\theta/C_0)s_0D|A|]$, proving the claim because
$s_0D=bD^{2/3}$. The empty-set assertion is immediate.
\end{proof}

The spatial choices have a fixed order: first $\alpha_0$ large and
$\ell$ small, then $\theta$ with $3\theta<\ell/2$, then $b$
small, and finally $D$ large. These estimates provide simultaneous
endpoint control for arbitrarily many prescribed cells. The next section
extends that control to every path visiting a cell during a short
imaginary-time interval.
